\scanpage{181}

\begin{center}
U N I V E R S I T E\\
DES SCIENCES ET TECHNIQUES\\
DU LANGUEDOC\\
\rule{2.8cm}{.4pt}\\
MA/FP
\end{center}
\uline{Diffusion} : Liste A\\
N° 3109

\begin{center}
\uline{AVIS DE SOUTENANCE D'UNE THESE DE DOCTEUR DE}\\
\uline{3ème CYCLE : TECTONIQUE-GEOPHYSIQUE- GEOCHIMIE}
\end{center}

\uline{Nom et Prénom du candidat} : RISPOLI René

\uline{Titre de la thèse} : Microtectonique et champ de contraintes dans les calcaires fins du Languedoc. Exemple des Matelles et du Cirque de Navacelles.

\uline{Date, heure et lieu de soutenance prévus} : 9 Juin 1981, 10h 30, Salle de Cours n° 6\\
U.S.T.L.

\uline{Nom du Directeur de thèse} : MM. les Professeurs MATTAUER, PROUST.

\uline{Désignation du laboratoire où la thèse est préparée} : Laboratoire de Géologie Structurale.

\uline{Analyse de l'ouvrage} : Aux Matelles (20 Km) au Nord de Montpellier, l'observation des microstructures tectoniques pyrénéennes dans les calcaires permet de décrire la géométrie du champ des contraintes et plus particulièrement les concentrations et déviations de $\sigma_1$, $\sigma_2$, $\sigma_3$. On propose à partir de ces données des modèles évolutifs pour les fentes, stylolites et failles. Ces microstructures s'associent entre elles dans les systèmes en échelon qui sont des cisaillements potentiels. Leur description fournit des renseignements sur leur fonctionnement mécanique.

Sur le site du Cirques de Navacelles (Causse du Larzac) la déformation tectonique pyrénéenne est observée. Deux phases de distension mésozoïques sont décrites. L'utilisation de l'analyse numérique des microfailles permet de définir la variation de $\mathrm R=\sigma_2-\sigma_3\,/\,\sigma_1-\sigma_3$ avec la profondeur z. L'utilisation des modèles de contraintes dans la croûte permet d'évaluer la poussée tectonique pyrénéenne entre 100 et 200 bars.

\uline{Durée approximative des recherches ayant abouti à la thèse} : 3 ans.

\uline{Composition du jury proposé} :

\begin{flushleft}
M. MATTAUER M\\
M. PROUST F.\\
M. ARTHAUD F\\
M. ETCHECOPAR A.\\
M. VASSEUR G.
\end{flushleft}

\begin{flushright}
p. Le Secrétaire Général\\
L'Attachée d'Administration Universitaire,\\
\uline{\unsure{Gentil}}\\
A. GENTIL.
\end{flushright}

\margin{181}
